
\documentclass[12pt]{article}

% Use algorithm and algpseudocode packages so \State and related commands are available

\newcommand{\duedate}{09/09/2025}
\newcommand{\assignment}{Problem Set 0} % Change to "Problem Set X"

% Change the following to your name and UNI.
\newcommand{\name}{Aksel Kretsinger-Walters, adk2164}
\newcommand{\email}{adk2164@columbia.edu}

% NOTE: Defining collaborators is optional; to not list collaborators, comment out the
% line below. Maximum of two collaborators per problem set.
%\newcommand{\collaborators}{Vig Vigerton (\texttt{UNI}), Alice Bob (\texttt{UNI})}
% No collaborators on PS 0

\makeatletter
\def\input@path{{../}{../../}{../../../}} % add as many parents as you need
\makeatother
\input{pset_template_RL.tex} %% DO NOT CHANGE THIS LINE

\begin{document}
\psetheader %% DO NOT CHANGE THIS LINE

\section*{Lecture 4 Summary: Policy Gradient Methods}

\subsection*{1. Motivation}

Value-based methods (e.g., Value Iteration, Q-Learning) learn value functions and then act greedily.
However, they can become unstable or inefficient with nonlinear function approximation.
Policy-gradient methods instead \emph{directly} parameterize and optimize the policy itself:
\[
		J(\theta) = \mathbb{E}_{\pi_\theta}\!\left[ \sum_{t=0}^{\infty} \gamma^t r_t \right].
\]
The goal is to find parameters $\theta$ that maximize $J(\theta)$.

\subsection*{2. Stochastic Policy Parameterizations}

\begin{table}[h!]
		\centering
		\footnotesize
		\begin{tabular}{@{}p{2.8cm}p{5.3cm}p{6.5cm}@{}}
				\hline
				\textbf{Setting} & \textbf{Example} & \textbf{Log-gradient} \\
				\hline
				\textbf{Discrete actions} &
				$\pi_\theta(a\mid s)=\dfrac{e^{f_\theta(s,a)}}{\sum_{a'} e^{f_\theta(s,a')}}$ &
				$\nabla_\theta \log \pi_\theta(a\mid s)=\nabla_\theta f_\theta(s,a)
				-\sum_{a'}\pi_\theta(a'\mid s)\nabla_\theta f_\theta(s,a')$ \\[0.8em]
				\textbf{Continuous actions} &
				Gaussian $\mathcal{N}(\mu_\theta(s),\sigma^2)$ &
				$\nabla_\theta \log \pi_\theta(a\mid s)=
				\dfrac{a-\mu_\theta(s)}{\sigma^2}\nabla_\theta\mu_\theta(s)$ \\
				\hline
		\end{tabular}
		\caption{Common stochastic policy parameterizations and their log-gradients.}
\end{table}

These forms make $\nabla_\theta \log \pi_\theta(a\mid s)$ directly computable by autodiff.

\subsection*{3. Policy Gradient Theorem}

The \emph{Policy Gradient Theorem} gives a simple, model-free expression for the gradient:
\begin{equation}
		\begin{aligned}
				\nabla_\theta J(\theta)
				&= \mathbb{E}_{s\sim d^{\pi_\theta},\,a\sim\pi_\theta}
				\big[\,Q^{\pi_\theta}(s,a)\,\nabla_\theta\log\pi_\theta(a\mid s)\,\big].
		\end{aligned}
\end{equation}
Key insight: the derivative does not include $\nabla_\theta d^{\pi_\theta}(s)$,
so we can estimate this gradient using on-policy samples.

\subsection*{4. REINFORCE Algorithm (Williams, 1992)}

A Monte Carlo estimator of the gradient:
\[
		\hat{g} = \sum_{t=0}^{T-1} R_t \, \nabla_\theta \log \pi_\theta(a_t\mid s_t),
		\quad
		R_t = \sum_{t' \ge t} \gamma^{t'-t} r_{t'}.
\]
% Corrected algorithm environment for REINFORCE
\begin{algorithm}[H]
		\small
		\caption{REINFORCE}
		\begin{algorithmic}[1]
    \State Initialize policy parameters $\theta$
    \For{each episode}
				\State Sample trajectory $\tau = (s_0, a_0, r_0, \ldots, s_T)$ under $\pi_\theta$
				\For{each timestep $t$ in $\tau$}
				\State Compute $R_t = \sum_{t' \ge t} \gamma^{t'-t} r_{t'}$
				\State $\theta \leftarrow \theta + \alpha R_t \nabla_\theta \log \pi_\theta(a_t \mid s_t)$
				\EndFor
    \EndFor
		\end{algorithmic}
\end{algorithm}

REINFORCE is unbiased but has high variance.

\subsection*{5. Variance Reduction via Baselines}

Subtracting a baseline $b(s)$ that does not depend on $a$ leaves the gradient unbiased:
\[
		\nabla_\theta J(\theta)
		= \mathbb{E}\!\left[ (Q^\pi(s,a) - b(s)) \, \nabla_\theta \log \pi_\theta(a\mid s) \right].
\]
Choosing $b(s) = V^\pi(s)$ minimizes variance and defines the \emph{Advantage}:
\[
		A^\pi(s,a) = Q^\pi(s,a) - V^\pi(s).
\]
The update then becomes:
\[
		\theta \leftarrow \theta + \alpha A_t \, \nabla_\theta \log \pi_\theta(a_t\mid s_t).
\]

\subsection*{6. Actor--Critic Methods}

In practice, $Q^\pi$ or $V^\pi$ are unknown and learned with a \emph{critic}.
The critic minimizes TD-error:
\[
		\delta_t = r_t + \gamma V_\omega(s_{t+1}) - V_\omega(s_t),
		\qquad
		\omega \leftarrow \omega - \beta \, \delta_t \nabla_\omega V_\omega(s_t).
\]
The actor then updates using $\delta_t$ as a low-variance estimate of the advantage:
\[
		\theta \leftarrow \theta + \alpha \, \delta_t \, \nabla_\theta \log \pi_\theta(a_t\mid s_t).
\]

\subsection*{7. Average Reward (Undiscounted) Case}

For continuing tasks, the objective is the average steady-state reward:
\[
		\rho^\pi = \sum_s d^\pi(s) \sum_a \pi(a\mid s) R(s,a),
\]
and the gradient takes the same form:
\[
		\nabla_\theta \rho^\pi
		= \mathbb{E}_{s,a \sim \pi}\!\left[
				A^\pi(s,a)\,\nabla_\theta \log \pi_\theta(a\mid s)
		\right].
\]

\subsection*{8. Conservative and Trust-Region Policy Updates}

Large policy changes can degrade performance.
To guarantee monotonic improvement:
\[
		\pi_{k+1} = (1-\alpha)\pi_k + \alpha \pi^*_k,
\]
or impose a KL constraint (as in TRPO):
\begin{equation}
		\begin{aligned}
				\max_{\theta}\;\;&
				\mathbb{E}\!\left[
						\frac{\pi_\theta(a\mid s)}{\pi_{\theta_{\mathrm{old}}}(a\mid
						s)}A^{\pi_{\theta_{\mathrm{old}}}}(s,a)
				\right]\\
				\text{s.t.}\;\;&
				\mathbb{E}\!\left[\mathrm{KL}\bigl(\pi_{\theta_{\mathrm{old}}}\,\|\,\pi_\theta\bigr)\right]\le\delta.
		\end{aligned}
\end{equation}
The practical approximation is PPO’s \emph{clipped surrogate} objective.

\subsection*{9. Key Takeaways}

\begin{itemize}
		\item Policy Gradient Theorem enables model-free differentiation of expected return.
		\item REINFORCE is unbiased but high-variance; baselines reduce variance.
		\item Actor--Critic uses bootstrapped value estimates to stabilize learning.
		\item TRPO/PPO constrain policy updates for reliable improvement.
		\item Policy-based methods handle continuous actions and stochasticity naturally.
\end{itemize}

\noindent
\textbf{In one line:}
Policy gradient methods learn by differentiating through the policy itself;
REINFORCE provides the estimator, baselines and critics control variance,
and trust-region methods make deep RL stable.

\end{document}
